% from ega4-9-fr.tex
\oldpage[IV-3]{86}

Il résulte de (4.2.7) et (4.2.8) que la condition (9.2.1, (i)) est satisfaite
pour les propriétés dont nous voulons démontrer qu'elles sont constructibles.
On est donc ramené au cas où $S$ est affine, noethérien et intègre, et à prouver
que $E$ et $E'$ sont constantes au voisinage du point générique $\eta$ de $S$~;
mais cela résulte de (9.8.3) et (9.5.5).

\begin{proposition}[9.8.6]
\label{IV.9.8.6-fr}
Avec les hypothèses et notations de (9.8.2), soit $i\in I$ tel que $(Z_i)_\eta$
soit maximal (autrement dit, une composante irréductible de
$\operatorname{Supp}(\mathcal F_\eta)$). Alors, il existe un voisinage $U$ de
$\eta$ dans $S$ tel que pour tout $s\in U$ et tout $\alpha\in J_{s,i}$, la
longueur géométrique de $\mathcal F_s$ en $x_{is\alpha}$ (relative à
$\boldsymbol{k}(s)$) (4.7.5) est égale à la longueur géométrique de
$\mathcal F_\eta$ en $x_i$ (relative à $\boldsymbol{k}(\eta)$).
\end{proposition}

On peut évidemment se borner au cas où $S=\operatorname{Spec}(A)$ est affine~;
montrons d'abord qu'on peut se ramener au cas où le sous-préschéma $(Z_i)_\eta$
de $X_\eta$, qui est réduit, est \emph{géométriquement intègre}. Il y a en effet
une extension finie $K'$ de $K=\boldsymbol{k}(\eta)$ telle que
$(((Z_i)_\eta)_{(K')})_{\mathrm{red}}$ soit géométriquement réduit et que les
composantes irréductibles de $((Z_i)_\eta)_{(K')}$ soient géométriquement
irréductibles (4.6.8). Procédons comme dans la démonstration de (9.7.7) en
considérant une sous-$A$-algèbre $A'$ de $K'$ ayant $K'$ pour corps des
fractions et finie sur $A$. On pose $S'=\operatorname{Spec}(A')$ et on considère
le morphisme fini surjectif $g:S'\to S$~; soient ensuite $X'=X_{(S')}$ et
$\mathcal F'=\mathcal F\otimes_{\mathcal O_S}\mathcal O_{S'}$, et soit $\eta'$
le point générique de $S'$. Pour tout $s'\in S'$, soit $s=g(s')$~; si $T$ est
une composante irréductible de $\operatorname{Supp}(\mathcal F_s)$, les
composantes irréductibles $T'_j$ de $T_{(\boldsymbol{k}(s'))}$ sont des
composantes irréductibles de $\operatorname{Supp}(\mathcal F'_{s'})$ et dominent
$T$ (4.2.7) et les multiplicités radicielles de $T$ par rapport à $\mathcal F_s$
et de chaque $T'_j$ par rapport à $\mathcal F'_{s'}$ sont les mêmes (4.7.9). Le
raisonnement de la première partie de (9.7.7) montre donc qu'on peut se borner à
démontrer la proposition pour $X'$ et $\mathcal F'$~; et en vertu du choix de
$K'$, les sous-préschémas réduits ayant pour espaces sous-jacents les composantes
irréductibles de $((Z_i)_\eta)_{(K')}$ sont géométriquement intègres (4.6.1).

Supposons donc désormais $(Z_i)_\eta$ géométriquement intègre~; alors (9.7.7),
il en est de même de $(Z_i)_s$ pour $s$ voisin de $\eta$~; la définition (4.7.5)
montre qu'il suffira donc de prouver que la \emph{longueur} du
$\mathcal O_{X_\eta,x_i}$-module $(\mathcal F_\eta)_{x_i}$ est égale à celle du
$\mathcal O_{X_s,x_{is}}$-module $(\mathcal F_s)_{x_{is}}$ (on a ici supprimé
l'indice $\alpha$, inutile par hypothèse). La question étant évidemment locale
sur $X$, on peut supposer (en se restreignant à un voisinage de $x_i$) que
$\mathcal F=\mathcal F_i$, donc que $\mathcal F_\eta$ est \emph{irredondant} sur
$X=\operatorname{Spec}(B)$ affine, et on écrira $Z$ au lieu de $Z_i$, et $x$
pour le point générique de $Z$ (et de $Z_\eta$). L'anneau noethérien $B$
contient donc $A$ comme sous-anneau, et $\mathcal F=\widetilde M$, où $M$ est un
$B$-module de type fini~; en outre, si $K$ est le corps des fractions de $A$,
le $B_{(K)}$-module $M_{(K)}$ est $\neq0$ et irredondant. Soit $q$ l'unique
élément de $\operatorname{Ass}(M_{(K)})$ et soit $\mathfrak p$ l'idéal premier
de $B$, image réciproque de $q$. En utilisant (5.11.1.1) comme dans la
démonstration de (9.7.6), on se ramène au cas où $B$ est \emph{intègre} et $M$
un sous-module $\neq0$ de $B$~; alors $\mathcal F_\eta$ est un
sous-$\mathcal O_{X_\eta}$-Module non nul de $\mathcal O_{Z_\eta}$, et en vertu
de (9.4.5), pour $s$ voisin de $\eta$, $\mathcal F_s$ est isomorphe à un
sous-$\mathcal O_{X_s}$-Module non nul de $\mathcal O_{Z_s}$~; comme $Z_s$ est
géométriquement intègre, les longueurs de $(\mathcal F_\eta)_x$ et de
$(\mathcal F_s)_{x_s}$ sont toutes deux égales à $1$, ce qui achève la
démonstration.

